\scanpage{32}
\textbf{C. R. Acad. Sc. Paris, t. 278 (27 mai 1974)}\hfill\textbf{Série A ---}\mspage{1403}\par\noindent\rule{\linewidth}{.35pt}\par

Par la bijection \(\psi_f\) (4.1), on en déduit une partition \((\mathrm{E}_i)_{i\in\mathrm{I}}\) de \(\mathrm{C}(\mathrm{K},\omega(f))\) et à chaque \(\mathrm{E}_i\) est associé un entier \(n_i\geqq0\). Alors \(\omega(f)\) et la famille \(((\mathrm{E}_i,n_i))_{i\in\mathrm{I}}\) sont des invariants, c'est-à-dire sont inchangés lorsqu'on remplace \(f\) par un plongement équivalent. On les appellera \textit{invariants} du plongement \(f\).

5.2. \textsc{Théorème de classification.} --- \textit{Deux plongements \(f\) et \(f'\) sont équivalents si et seulement s'ils ont mêmes invariants. De plus, tout système d'invariants sur \(\mathrm{K}\) : \(\omega,\ (\mathrm{E}_i,n_i)_{i\in\mathrm{I}}\) provient d'une classe de plongements.}

{\small (*) Séance du 17 avril 1974.

(\textsuperscript{1}) J. R. \textsc{Edmonds}, \textit{A combinatorial representation for polyhedral surfaces} (\textit{Not. A. M. S.}, 7, 1960, p. 646).

(\textsuperscript{2}) J. F. P. \textsc{Hudson}, \textit{Piece-Wise Linear Topology}, W. A. Benjamin Inc., 1969.

(\textsuperscript{3}) R. L. \textsc{Wilder}, \textit{Topology of Manifolds} (\textit{A. M. S.}, Coll. Publ., 32, 1949).

(\textsuperscript{4}) E. C. \textsc{Zeeman}, \textit{Seminar on Combinatorial topology}, I. H. E. S., 1963, Exp. I à VIII.}

\begin{flushright}
\textit{Université des Sciences et Techniques}\\
\textit{du Languedoc,}\\
\textit{Mathématiques,}\\
\textit{Place Eugène-Bataillon,}\\
34000 \textit{Montpellier.}
\end{flushright}
% archive sheet number: 31

